% ============================================================================
% 1. Introduction  +  2. Related Work
% ============================================================================
\section{Introduction}\label{sec:intro}

Self-attention is the computational core of the Transformer architecture
\cite{vaswani2017}. In its standard form it computes, for every pair of positions,
a dot product between a query and a key vector, divides that logit by
$\sqrt{\dk}$, applies a softmax normalisation, and mixes value vectors with the
resulting weights. The $1/\sqrt{\dk}$ factor is justified by a variance argument:
if the components of $d_k$-dimensional queries and keys are treated as independent
random variables with unit variance, the dot product has variance $d_k$, so
dividing by $\sqrt{\dk}$ keeps the logit variance constant and prevents the
softmax from saturating prematurely \cite{vaswani2017}. This rationale has been
quoted across hundreds of follow-up works, yet the scaling constant itself is
seldom the object of a controlled empirical test: it is inherited, not chosen.

Alternatives are mathematically trivial to implement. Dividing by $\dk$ directly
removes the square root; introducing an explicit softmax temperature
$\tau$ \cite{hinton2015} rescales the same logits by a tunable constant. Both
leave the parameter count, the computational graph, and the memory traffic
unchanged, and both plausibly alter training dynamics by reshaping the entropy of
the attention distribution \cite{zhang2022temp,zhai2023,dong2021}. Recent
architectural work has begun to treat the scaling function as a first-class
design axis~--- proposing input-size-aware softmax rescaling \cite{nakanishi2025},
softplus-based alternatives \cite{gao2025}, and temperature adjustments for
length extrapolation \cite{chi2023}~--- but the empirical base remains thin:
most of these studies compare architectural systems end-to-end, where many
factors change at once and a single seed per configuration is common.

Whether a different scaling constant helps or hurt is therefore an empirical
question, and answering it credibly requires more than running two configurations
once. Small models trained for short budgets exhibit substantial seed-to-seed
variability \cite{bouthillier2021,picard2021}, so single-run comparisons mostly
measure noise \cite{dodge2019}. Matched-pair designs, in which competing variants
share the same random initialization and the same data order, are a more
sensitive instrument: they remove variance sources that a between-run comparison
cannot control, and they support paired statistical tests on per-seed differences
\cite{dror2018}.

This study reports a comparison of exactly that kind. We implemented a compact
pre-LN Transformer encoder \cite{xiong2020,ba2016} whose attention module admits
three scaling functions, verified by unit tests that the variants share
bit-identical initialization and produce correctly masked computations, and then
trained fixed grids of 3 variants $\times$ 25 matched seeds on four text datasets
under frozen hyperparameter budgets. Four research questions structure the
investigation:

\begin{itemize}[leftmargin=1.8em]
  \item \textbf{RQ1.} Does standard scaled dot-product attention outperform the
        alternative scaling functions on predictive accuracy under a matched
        budget?
  \item \textbf{RQ2.} Do the variants differ in computational efficiency
        (training time, inference latency, memory) when the scaling change is
        parameter-free and adds no arithmetic?
  \item \textbf{RQ3.} How do attention entropy, head diversity, and gradient
        magnitudes respond to the scaling factor, and do those responses explain
        the performance outcomes?
  \item \textbf{RQ4.} Are the observed performance differences statistically
        significant under paired tests with multiple-comparison correction?
\end{itemize}

The preregistered hypothesis H1~--- that at least one alternative scaling
function would significantly outperform standard attention on a relevant metric
--- is \emph{partially falsified} at the Stage-1 operating point: standard
attention significantly outperformed both alternatives on AG News, the
alternatives were not significantly worse on SST-2 or IMDB, and no variant
separated from the rest on WikiText-2 with sufficient statistical power. This
revision completes the programme and sharpens the picture further: with the
corrected causal mask the WikiText-2 comparison is properly powered (25 seeds
per arm) and shows no divisor effect; the head-dimension grid shows no
amplification with $\dk$; and at a documented replica operating point the
Stage-1 AG News effect does not reproduce, bounding its generality. The
variants were exactly equivalent in measured computational cost, so the
accuracy deficit (where present) is not offset by any efficiency gain.

The contribution of this paper is sixfold. (i) A clean, parameter-free
experimental manipulation~--- a single divisor in one line of code~--- evaluated
under a seed-matched paired design across 787 completed runs in three stages
(GPU scale, replica completion, and an operating-point programme), with paired
$t$-tests, Wilcoxon signed-rank tests, Holm--Bonferroni
correction, and effect sizes. (ii) Evidence that the standard $\sqrt{\dk}$
divisor's advantage is conditional: significant on a 4-class
topic-classification task at full data scale and a short budget, absent on
binary sentiment tasks, absent in the replica completion of the core grid, and
not offset by any computational trade-off. (iii) A corrected causal-mask
language-modelling protocol with formal mask-verification tests and leak probes
that explain~--- rather than merely discard~--- the artifactually low
perplexities of the original runs. (iv) Mechanistic diagnostics linking the
scaling factor to attention entropy, inter-head similarity, and gradient
magnitudes, with the measured $1/s$ gradient scaling matching theory, and a
replica-scale regime analysis (raw logit scale versus divisor) that explains
when the divisor can and cannot matter. (v) A quantitative operating-point
programme (Stage 3): a Monte-Carlo-validated power law for the divergence
between divisors as a function of raw logit variance, a logit-gain sweep and
regime-restored training runs that cross the bite threshold, frontier
stabilisers (QK-Norm, soft-capping) tested in context, and an
inference-time divisor swap in pretrained GPT-2 decoders showing the contrast
is inert at replica scale but worth a factor of 94 in perplexity at pretrained
scale, with a GPU script extending the identical protocol to 1B--7B models.
(vi) A fully reproducible package---
configurations, unit tests, analysis scripts, per-seed records, and this
manuscript~--- released under an MIT licence.\footnote{\url{https://github.com/razahasan2000/Attention-Scaling}}

The remainder of the paper is organised as follows.
Section~\ref{sec:related} situates the study in the literature on attention
scaling, attention dynamics, and experimental methodology.
Section~\ref{sec:method} describes the variants, architecture, datasets,
matched-seed design, and statistical analysis.
Section~\ref{sec:results} reports the outcomes of all completed grids.
Section~\ref{sec:discussion} interprets the findings, and
Sections~\ref{sec:limitations} and \ref{sec:conclusion} state limitations and
conclusions.

% ============================================================================
\section{Related Work}\label{sec:related}

\subsection{Attention scaling and softmax temperature}
The scaled dot-product formulation with the $\sqrt{\dk}$ divisor was introduced
with the Transformer \cite{vaswani2017} and inherited by essentially every
successor, including pre-LN variants that reorganised the residual stream for
stability \cite{xiong2020,ba2016}. The temperature of a softmax~--- the constant
that divides the logits before exponentiation~--- entered deep learning through
knowledge distillation, where Hinton et al. \cite{hinton2015} showed that
elevated temperatures expose more of the teacher's class-similarity structure.
In attention, Zhang et al. \cite{zhang2022temp} demonstrated that a tuned
attention temperature matters for summarization distillation, and subsequent
work has made the temperature task- or layer-dependent. Nakanishi \cite{nakanishi2025}
proposed Scalable-Softmax, which rescales logits as a function of the input
vector size to counteract attention ``fading'' over long contexts, and Gao et
al. \cite{gao2025} replaced softmax with a length-scaled softplus for similar
reasons. On the theoretical side, Dong et al. \cite{dong2021} showed that pure
attention networks lose rank doubly exponentially with depth~--- a collapse that
residual connections and, indirectly, the sharpness of the attention
distribution modulate~--- while Chiang and Cholak \cite{chiang2022} analysed how
softmax attention can plateau in probability space and how the scaling interacts
with easy fixes. Zhai et al. \cite{zhai2023} connected training stability to
attention entropy, showing that entropy collapse (near-uniform attention) can be
prevented by reparameterisation, which positions the scaling constant as a
direct lever on the same quantity. What is missing from this literature is a
controlled, seed-matched measurement of how much the classical constant
$\sqrt{\dk}$ itself matters on standard NLP benchmarks~--- the gap this study
addresses.

\subsection{Attention heads, diversity, and their measurement}
Multi-head attention analyses show that individual heads specialise, and that
many heads can be pruned with little loss \cite{voita2019,michel2019}. Head
specialisation is precisely the property that a flatter attention distribution
should erode: if all heads attend near-uniformly, their outputs become similar.
We therefore measure inter-head cosine similarity of attention distributions and
mean attention entropy as mechanistic mediators between the scaling constant and
task performance (Section~\ref{sec:diagnostics}). Related measurement tools ---
attention-map similarity and pruning sensitivity~--- have a long history, but
they have not previously been applied to isolate the effect of the logit scaling
constant under matched training conditions.

\subsection{Positional and architectural context}
Scaling interacts with other attention modifiers studied for length
generalisation: additive attention biases (ALiBi) \cite{press2022}, rotary
rescaling methods such as YaRN \cite{peng2024}, and positional-encoding choices
\cite{kazemnejad2023,chi2023}. These works modify attention for long-context
capability, whereas the present study deliberately freezes the architecture ---
including a fixed maximum sequence length of 128~--- and varies only the scalar
divisor. The two lines are complementary: a reliable answer about the default
divisor is a precondition for reasoning about more exotic rescalings.

\subsection{Experimental methodology: variance and significance}
Deep-learning benchmarks are dominated by variance across random seeds,
compute-minority reports, and incomplete hyperparameter searches
\cite{bouthillier2021,dodge2019}. In NLP specifically, Dror et al. \cite{dror2018}
found that significance testing is often ignored or misused, and recommended
matched designs with paired tests where possible. Even seed sensitivity alone
can flip conclusions in computer-vision benchmarks \cite{picard2021}. The
methodology of this study follows these recommendations directly: 25 matched
seeds per variant, within-seed paired differences as the unit of analysis,
paired $t$-tests with Wilcoxon signed-rank confirmation \cite{wilcoxon1945},
Holm--Bonferroni family-wise correction across the two comparisons per dataset
\cite{holm1979}, and paired effect sizes $d_z$ with conventional magnitude
labels \cite{cohen1988}. To our knowledge this is the first application of a
fully seed-matched paired protocol to the question of attention logit scaling.
